% open-logic.tex
% compiles to produce complete text using open-logic-dev style

\documentclass[../include/open-logic-part]{subfiles}

\begin{document}

\clearpage

\begin{editorial}
এই ফাইল ওপেন লজিক প্রকল্পের সব অন্তর্ভুক্ত বিষয়বস্তু লোড করে।
এই ধরনের সম্পাদকীয় টীকা দেখা গেলে বুঝতে হবে, বিষয়বস্তু কীভাবে
উপস্থাপিত হবে তা বিবেচনা না করেই ফাইলটি সংকলিত হয়েছে।
এটি পড়তে পেলে সেই PDF থেকে পড়ানো বা পড়াশোনা করা সম্ভবত
\emph{সমীচীন নয়}।

শিক্ষাদান বা স্বশিক্ষার জন্য আরও উপযুক্ত পাঠ তৈরির বহু ব্যবস্থা
ওপেন লজিক প্রকল্পে আছে। যেমন, সাধারণ বিন্যাসে সব যৌক্তিক অপারেটরকে
মৌলিক ধরা হয় এবং প্রমাণে সব অপারেটরের সব ক্ষেত্র আলোচনা করা হয়।
কিন্তু তার কিছু ক্ষেত্র অনুশীলন হিসেবে রাখা অনেক ভাল।
ওপেন লজিক প্রকল্পের কাজও চলছে। সহযোগিতা ও উন্নতিতে উৎসাহ দেওয়ার
জন্য খসড়া বিষয়বস্তু, অনুশীলনহীন বিভাগ ইত্যাদিও অন্তর্ভুক্ত হয়।
কোনো পাঠক্রমের জন্য তৈরি PDF-এ এই বিভাগগুলি বাদ দেওয়া হবে।

শিক্ষাদান ও পড়াশোনার জন্য আরও উপযুক্ত PDF পেতে
\href{http://builds.openlogicproject.org/}{ওপেন লজিক প্রকল্পের
ওয়েবসাইটে পাওয়া নমুনা পাঠক্রমগুলি} দেখো। নিজের সংস্করণ তৈরির
শুরু হিসেবে
\href{https://github.com/OpenLogicProject/OpenLogic/tree/master/courses/sample}{নমুনা
নির্দেশফাইল} ব্যবহার করতে পারো, অথবা আরও উন্নত উদাহরণের জন্য
এই প্রকল্প থেকে তৈরি অন্যান্য পাঠ্যবইয়ের উৎস দেখো।
\end{editorial}

\olimport[sets-functions-relations]{sets-functions-relations-complete}

\olimport[propositional-logic]{propositional-logic}

\olimport[first-order-logic]{first-order-logic}

\olimport[model-theory]{model-theory}

\olimport[computability]{computability}

\olimport[turing-machines]{turing-machines}

\olimport[incompleteness]{incompleteness}

\olimport[second-order-logic]{second-order-logic}

\olimport[lambda-calculus]{lambda-calculus}

\olimport[many-valued-logic]{many-valued-logic}

\olimport[normal-modal-logic]{normal-modal-logic}

\olimport[applied-modal-logic]{applied-modal-logic}

\olimport[intuitionistic-logic]{intuitionistic-logic}

\olimport[counterfactuals]{counterfactuals}

\olimport[set-theory]{set-theory}

\olimport[methods]{methods}

\olimport[history]{history}

\olimport[reference]{reference}

\end{document}
